% Einheit 1 von 5 – Bernoulli-Experiment und Bernoulli-Kette (bank/binomialverteilung/e1.jsonl, zusammenbau v0.1)
%% TODO Zeitmarke Einheit 1: Marken-Zeile nicht lesbar
%% TODO Zweigzeile Teil 1: was hier gelernt wird (Satz in Schülersprache aus dem Einheitstitel) – Einheit 1
\einheitenkopf[e1]{Einheit 1 von 5 · Bernoulli-Experiment und Bernoulli-Kette}
\zweigzeile{Abitur GK}
\setcounter{aufgabe}{10}

%% TODO Titel: Ich-kann-Satz fehlt in der Bank (Platzhalter: Kettenname) – Nr. 11
%% TODO Anweisung: ein Satz über der ersten Teilaufgabe fehlt in der Bank – Nr. 11
\begin{aufgabe}{Modell erkennen}
%% TODO \rechenplatz{n}: Schrittzahl des Grundfalls fehlt in der Bank (nur bei fester Schreibform, 2.3 b)
\begin{teile}
\teil Eine Münze wird 20-mal geworfen; gezählt wird, wie oft Wappen fällt. Bernoulli-Kette oder nicht? Nenne bei „nein“ die verletzte Bedingung. \kreuz{Bernoulli-Kette} \\ \kreuz{keine Bernoulli-Kette}
\teil Eine Münze wird einmal geworfen; Treffer ist „Wappen“. Mit welcher Wahrscheinlichkeit tritt der Treffer ein? p = \leerfeld
\teil Ein Glücksrad hat 5 gleich große Felder, 2 davon sind blau. Es wird 12-mal gedreht; X ist die Anzahl der Drehungen mit Blau. Gib die Anzahl der Versuche n und die Trefferwahrscheinlichkeit p an und schreibe die Verteilung von X kurz. X ~ B( \leerfeld ; \leerfeld )
\teil Ein Glücksrad bleibt mit der Wahrscheinlichkeit $0{,}15$ auf Gold stehen. Es wird zweimal gedreht. Genau einmal Gold – mit welcher Wahrscheinlichkeit? \leerfeld
\teil Ein Tetraeder mit den Zahlen 1 bis 4 wird dreimal geworfen. Bei 1, 2 oder 3 wird eine rote Perle aufgefädelt, sonst eine weiße. Zeige, dass mit der Wahrscheinlichkeit $\frac{27}{32}$ mindestens zwei rote Perlen auf der Schnur sind \hfill (Abitur 2023 LK).
\teil In einem Briefzentrum kommen 4\,\% der Briefe beschädigt an. 150 Briefe aus der Tagespost werden zufällig ausgewählt; untersucht wird die Anzahl der beschädigten. Begründe, dass die Binomialverteilung verwendet werden kann \hfill (Abitur 2022 GK).
\teil Von 30 geprüften Laptops sind 5 defekt. Nacheinander werden 8 Laptops zufällig ausgewählt, ohne sie zurückzulegen. Beurteile, ob die Anzahl der defekten Laptops unter den ausgewählten binomialverteilt ist \hfill (Abitur 2018 GK).
\teil Ein Glücksrad hat zehn gleich große Felder, sieben mit „A“ und drei mit „B“. Je Spiel wird zweimal gedreht: „BB“ bringt den Hauptpreis, „AA“ einen Trostpreis. Beurteile die Aussagen (1) „Eine Drehung, bei der nur der Buchstabe betrachtet wird, ist ein Bernoulli-Experiment.“ und (2) „Werden mehrere Spiele gespielt und jeweils festgehalten, ob Trost- oder Hauptpreis vergeben wird, liegt eine Bernoulli-Kette vor.“ \hfill (Abitur 2018 GK).
\end{teile}
\end{aufgabe}

%% TODO Titel: Ich-kann-Satz fehlt in der Bank (Platzhalter: Kettenname) – Nr. 12
%% TODO Anweisung: ein Satz über der ersten Teilaufgabe fehlt in der Bank – Nr. 12
\begin{aufgabe}{Modell erkennen – Fehler finden, Begründen, Anwendung, Darstellungswechsel}
\begin{teile}
\teil Mia sagt: „Aus einer Kiste mit 15 Äpfeln, von denen 4 faul sind, nimmt Jan 5 Äpfel heraus. Jeder Apfel ist faul oder gut, also ist die Anzahl der faulen Äpfel binomialverteilt.“ Finde den Fehler.
\teil In Deutschland leben rund 84 Millionen Menschen. Für eine Umfrage werden 1\,000 davon zufällig ausgewählt, niemand wird zweimal befragt. Begründe, warum man trotzdem mit der Binomialverteilung rechnen darf.
\teil Ein Flugzeug hat zwei Triebwerke; jedes fällt auf einem Flug unabhängig vom anderen mit der Wahrscheinlichkeit $0{,}001$ aus. Das Flugzeug kann mit einem Triebwerk sicher landen. Wie wahrscheinlich ist eine sichere Landung, also mindestens ein laufendes Triebwerk?
\teil Ein Glücksrad bringt bei jeder Drehung mit der Wahrscheinlichkeit $0{,}25$ einen Gewinn (G), sonst nichts (V). Es wird zweimal gedreht. Zeichne das Baumdiagramm mit allen Wahrscheinlichkeiten.

\baumzwei{G/,V/}{G/,V/}{G/,V/}
\end{teile}
\end{aufgabe}
